% ============================================================
% PASO 1 — Entalpía molar rigurosa
%          Gas real con ecuación virial truncada
%          Herrmann et al. (2009), Hyland & Wexler (1983)
%
%   ψ_w  →  \vaporMoleFraction   (y_v)
%   1-ψ_w → \gasMoleFraction     (y_g)
%   T    →  \T
% ============================================================
\NewDocumentCommand{\eqEnthalpyRigorous}{s}{%
    \begin{equation}
        \IfBooleanTF{#1}{}{\label{eq:enthalpy:rigorous}}
        \hMolar\!\left(\T,\, \molarVolume,\, \vaporMoleFraction\right)
        =
        \hMolarRef
        + \gasMoleFraction\, \hMolarGas
        + \vaporMoleFraction\, \hMolarVapor
        + \universalConstantGas\T
          \left[
            \left(
              \constBm
              - \T\frac{\mathrm{d}\constBm}{\mathrm{d}\T}
            \right)
            \frac{1}{\molarVolume}
            +
            \left(
              \constCm
              - \frac{\T}{2}\frac{\mathrm{d}\constCm}{\mathrm{d}\T}
            \right)
            \frac{1}{\molarVolume^2}
          \right]
    \end{equation}
    \IfBooleanTF{#1}{}{%
        \noindent donde:
        \begin{itemize}
            \item $\hMolar$ es la entalpía molar de la mezcla gaseosa \unit{\frac{kJ}{mol}}.
            \item $\T$ es la temperatura seca de la mezcla \unit{K}.
            \item $\molarVolume$ es el volumen molar de la mezcla
                  \unit{\frac{m^3}{mol}}.
            \item $\vaporMoleFraction$ es la fracción molar del vapor
                  en la mezcla.
            \item $\hMolarRef$ es la entalpía molar en el estado de
                  referencia \unit{\frac{kJ}{mol}}.
            \item $\gasMoleFraction = 1 - \vaporMoleFraction$ es la
                  fracción molar del gas seco en la mezcla.
            \item $\hMolarGas$ es la entalpía molar estándar del gas
                  seco \unit{\frac{kJ}{mol}}.
            \item $\hMolarVapor$ es la entalpía molar estándar del
                  vapor \unit{\frac{kJ}{mol}}.
            \item $\universalConstantGas$ es la constante universal de los
                  gases ideales
                  ($\universalConstantGas = 8.314\,\unit{\frac{J}{mol\cdot K}}$).
            \item $\constBm$ es el segundo coeficiente virial de la
                  mezcla \unit{\frac{m^3}{mol}}; cuantifica las
                  interacciones entre pares de moléculas.
            \item $\constCm$ es el tercer coeficiente virial de la
                  mezcla \unit{\frac{m^6}{mol^2}}; cuantifica las
                  interacciones entre tríos de moléculas.
        \end{itemize}
    }%
}

% ============================================================
% PASO 2 — Gas ideal: términos viriales → 0
%          Hipótesis: z_v = z_g = 1  →  v̄ → ∞
% ============================================================
\NewDocumentCommand{\eqEnthalpyIdealMolar}{s}{%
    \begin{equation}
        \IfBooleanTF{#1}{}{\label{eq:enthalpy:ideal:molar}}
        \hMolar\!\left(\T,\, \vaporMoleFraction\right)
        \approx
        \hMolarRef
        + \gasMoleFraction\, \hMolarGas
        + \vaporMoleFraction\, \hMolarVapor
    \end{equation}
}

% ============================================================
% PASO 3 — Relaciones entre entalpías molares y propiedades
%          caloríficas + relación con Y
%
%   y_v = ψ_w = Y·M_g / (M_v + Y·M_g)
%   y_g = 1 - y_v
%
%   CORRECCIÓN: \latentHeat (= Q_L, calor latente total) fue
%   reemplazado por \vapLatentHeat (= λ, calor latente
%   específico de vaporización, kJ/kg).
%   CORRECCIÓN: \omega reemplazado por \humidityRatio (= Y).
% ============================================================
\NewDocumentCommand{\eqEnthalpyConversion}{s}{%
    \begin{align}
        \IfBooleanTF{#1}{}{\label{eq:enthalpy:conversion}}
        \hMolarGas
            &= \cpGas \cdot \left(\T - \Tref\right) \cdot \gasMolecularWeight
            \notag \\
        \hMolarVapor
            &= \left[\vapLatentHeat + \cpVapor \cdot \left(\T - \Tref\right) \right]
               \cdot \vaporMolecularWeight
            \notag
    \end{align}
    \IfBooleanTF{#1}{}{%
        \noindent donde:
        \begin{itemize}
            \item $\hMolarGas$ es la entalpía molar estándar del gas
                  seco, calculada a partir de su calor específico
                  \unit{\frac{kJ}{mol}}.
            \item $\hMolarVapor$ es la entalpía molar estándar del
                  vapor, que incluye la contribución del calor
                  latente de vaporización \unit{\frac{kJ}{mol}}.
            \item $\cpGas$ es el calor específico a presión constante
                  del gas seco \unit{\frac{kJ}{kg \cdot K}}.
            \item $\cpVapor$ es el calor específico a presión constante
                  del vapor \unit{\frac{kJ}{kg \cdot K}}.
            \item $\T$ es la temperatura seca de la mezcla \unit{K}.
            \item $\Tref$ es la temperatura de referencia 
                  del estado estándar, habitualmente 
                  $\Tref = \qty{273.15}{K}$ 
                  o $\qty{298.15}{K}$ según la convención adoptada.
            \item $\vapLatentHeat$ es el calor latente de
                  vaporización del componente que se condensa a la temperatura $\T$
                  \unit{\frac{kJ}{kg}}.
            % \item $\gasMolecularWeight$ es el peso molecular del gas seco \unit{\frac{kg}{mol}}.
            % \item $\vaporMolecularWeight$ es el peso molecular del vapor \unit{\frac{kg}{mol}}.
            % \item $\humidityRatio$ es la razón de humedad másica, definida como la masa de vapor por unidad de masa de gas seco ($\humidityRatio = \nicefrac{\vaporMass}{\gasMass}$) \unit{\frac{kg_{\vapor}}{kg_{\gas}}}.
            % \item $\vaporMoleFraction$ es la fracción molar del vapor
                  % en la mezcla, obtenida a partir de $\humidityRatio$,
                  % $\gasMolecularWeight$ y $\vaporMolecularWeight$.
            % \item $\gasMoleFraction$ es la fracción molar del gas seco
                  % en la mezcla.
        \end{itemize}
    }%
}

% ============================================================
% PASO 4 — Entalpía específica expandida
%          Por unidad de masa de gas no condensable
%
%   h = c_{p,g}·T + Y·(λ + c_{p,v}·T)
%
%   CORRECCIÓN: \latentHeat → \vapLatentHeat
%   CORRECCIÓN: \omega      → \humidityRatio (Y)
% ============================================================
\NewDocumentCommand{\eqEnthalpyExpanded}{s}{%
    \begin{equation}
        \IfBooleanTF{#1}{}{\label{eq:enthalpy:expanded}}
        \hSpec
        =
        \cpGas \cdot \left(\T - \Tref \right)
        + \humidityRatio \cdot 
          \!\left[\vapLatentHeat + \cpVapor \cdot \left(\T - \Tref \right)\right]
        \quad \text{expresado en }
        \left[\frac{\text{kJ}}{\text{kg}_{\text{\gas}}}\right]
    \end{equation}
    \IfBooleanTF{#1}{}{%
        \noindent donde:
        \begin{itemize}
            \item $\hSpec$ es la cantidad de energía cedida a la mezcla por kg de gas seco; es decir, es la entalpía específica de la mezcla gaseosa
                  por unidad de masa de gas seco \unit{\frac{kJ}{kg_{\gas}}}.
            \item $\cpGas$ es el calor específico a presión constante
                  del gas seco \unit{\frac{kJ}{kg \cdot K}}.
            \item $\T$ es la temperatura seca de la mezcla \unit{K}.
            \item $\humidityRatio$ es la razón de humedad de la mezcla
                  \unit{\frac{kg_{\vapor}}{kg_{\gas}}}.
            \item $\vapLatentHeat$ es el calor latente de
                  vaporización del componente que se condensa a la temperatura $\T$
                  \unit{\frac{kJ}{kg}}.
            \item $\cpVapor$ es el calor específico a presión constante
                  del vapor \unit{\frac{kJ}{kg \cdot K}}.
        \end{itemize}
    }%
}

\NewDocumentCommand{\eqEnthalpyExpandedDevelop}{s}{%
    \begin{equation*}
        \IfBooleanTF{#1}{}{\label{eq:enthalpy:expanded:develop}}
        \hSpec
        =
        \left(\cpGas + \humidityRatio \cdot \cpVapor \right)  \left(\T - \Tref \right)
        + \humidityRatio \cdot \vapLatentHeat,
    \end{equation*}
}


% ============================================================
% PASO 5 — Calor específico húmedo
%
%   c_{p,H} = c_{p,g} + Y·c_{p,v}
%
%   CORRECCIÓN: \omega → \humidityRatio (Y)
% ============================================================
\NewDocumentCommand{\eqCpHumid}{s}{%
    \begin{equation}
        \IfBooleanTF{#1}{}{\label{eq:enthalpy:cp:humid}}
        \cpHumid = \cpGas + \humidityRatio \cdot \cpVapor
    \end{equation}
    \IfBooleanTF{#1}{}{%
        \noindent donde:
        \begin{itemize}
            \item $\cpHumid$ es el calor específico húmedo de la
                  mezcla: cantidad de calor necesaria para elevar en
                  un kelvin la temperatura de la unidad de masa de gas
                  seco junto con el vapor que le acompaña
                  \unit{\frac{kJ}{kg_{\gas}}\cdot K }.
            \item $\cpGas$ es el calor específico a presión constante
                  del gas seco \unit{\frac{kJ}{kg \cdot K}}.
            \item $\humidityRatio$ es la razón de humedad de la mezcla
                  \unit{\frac{kg_{\vapor}}{kg_{\gas}}}.
            \item $\cpVapor$ es el calor específico a presión constante
                  del vapor \unit{\frac{kJ}{kg \cdot K}}.
        \end{itemize}
    }%
}

% ============================================================
% PASO 6 — Entalpía específica compacta (forma final clásica)
%
%   h = c_{p,H}·T + λ·Y
%
%   CORRECCIÓN: \latentHeat → \vapLatentHeat
%   CORRECCIÓN: \omega      → \humidityRatio (Y)
% ============================================================
\NewDocumentCommand{\eqEnthalpyCompact}{s}{%
    \begin{equation}
        \IfBooleanTF{#1}{}{\label{eq:enthalpy:compact}}
        \hSpec
        =
        \cpHumid \cdot \left(\T - \Tref \right) + \vapLatentHeat\ \cdot \humidityRatio
      %   \quad
      %   \left[\frac{\text{kJ}}{\text{kg}_{\text{\gas}}}\right]
    \end{equation}
    \IfBooleanTF{#1}{}{%
        \noindent donde:
        \begin{itemize}
            \item $\hSpec$ es la entalpía específica de la mezcla gaseosa
                  por unidad de masa de gas seco \unit{\frac{kJ}{kg_{\gas}}}.
            \item $\cpHumid$ es el calor específico húmedo de la
                  mezcla, definido en la
                  ecuación~\ref{eq:enthalpy:cp:humid}
                  \unit{\frac{kJ}{kg_{\gas}}\cdot K }.
            \item $\T$ es la temperatura seca de la mezcla \unit{K}.
            \item $\vapLatentHeat$ es el calor latente de
                  vaporización del componente que se condensa a la temperatura $\T$
                  \unit{\frac{kJ}{kg}}.
            \item $\humidityRatio$ es la razón de humedad de la mezcla
                  \unit{\frac{kg_{\vapor}}{kg_{\gas}}}.
        \end{itemize}
    }%
}